% =============================================================================
% SECTION 3: STRUCTURAL CLOSURE ASSESSMENT
% =============================================================================
% Target: ~3 pages PRD format (dense, table-heavy)
% This is the paper's PRIMARY NOVELTY.
% =============================================================================

\section{Structural Closure of the Minimal Framework}
\label{sec:closure}

The nonsingular bounce described in Sec.~\ref{sec:bounce} is mathematically
well-defined and singularity-free.  We now ask: does this framework produce
testable consequences beyond the bounce itself?  We systematically investigate
all logical routes from the ECH bounce to late-time observables, organizing
them into fifteen branches labeled A--O.  Each branch is pursued until it
either produces a distinctive, testable prediction or encounters a structural
barrier---a mechanism-independent obstruction rooted in symmetry, scaling, or
conservation law.

The outcome is a catalog of fourteen structural barriers spanning five failure
modes.  No first-principles route from the bounce to dark energy survives, and
no distinctive observational signature of the bounce is accessible to current
or planned experiments.  One indirect handle remains: spectator ALP
birefringence, treated in Sec.~\ref{sec:alp}.


% ---------------------------------------------------------------------------
\subsection{Methodology}
\label{sec:closure_method}

We classify candidate connections between the bounce and late-time physics into
three tiers:

\begin{enumerate}
\item \emph{Direct derivation} (Branches A--G): attempts to derive a late-time
  dark energy density $\rho_\mathrm{DE}$ or equation of state $w=-1$ from the
  ECH Lagrangian and bounce dynamics, without external input.

\item \emph{Observable signatures} (Branches H, K--M): attempts to identify
  distinctive features in the tensor spectrum, scalar perturbations,
  gravitational wave background, or relic abundances that could distinguish the
  ECH bounce from other cosmological scenarios.

\item \emph{State selection and vacuum determination} (Branches I, J, N, O):
  attempts to use the bounce as a mechanism for constraining dark-sector
  initial conditions, triggering phase transitions, or selecting a vacuum
  energy scale.
\end{enumerate}

Each branch is tested against explicit calculations.  A branch is closed when
a structural barrier is identified---that is, an obstruction that cannot be
removed by parameter adjustment within the minimal framework or its natural
extensions (PGT, dynamical Immirzi, environmental masses).  A branch is
classified as \emph{generic} when the result, though calculable, is shared by
all bouncing cosmologies and cannot distinguish ECH from alternatives.


% ---------------------------------------------------------------------------
\subsection{Fourteen Structural Barriers}
\label{sec:barriers}

Table~\ref{tab:barriers} catalogs the fourteen barriers encountered across
all fifteen branches.  We summarize each in turn.

\begin{table*}[t]
\centering
\caption{Fourteen structural barriers closing all routes from the ECH bounce to
  late-time observables.  Each barrier is labeled by the branch in which it
  first appears, together with its physical mechanism and failure mode
  classification.}
\label{tab:barriers}
\begin{tabular}{clllll}
\hline\hline
\# & Barrier name & Branch & Mechanism blocked & Root cause & Failure mode \\
\hline
1  & Mass-coupling lock
   & A & PGT torsion as DE
   & $m \propto |t_3|^{-1/2}$,\; $g \propto (M_\mathrm{Pl}\sqrt{|t_3|})^{-1}$
   & I \\
2  & Topological-shift duality
   & B & Nieh--Yan pseudoscalar
   & Shift symmetry $\Leftrightarrow$ geometric content exclusive
   & I \\
3  & Scalar-tensor universality
   & C & Environmental mass on FRW
   & $T_0 = Q_0 = 0$ on FRW;\; only $R(g)$ survives
   & I \\
4  & Planck suppression
   & D & Disformal/connection coupling
   & $1\;\partial\phi$/vertex;\; disformal needs 2;\; dim-8 $\sim 10^{-122}$
   & I \\
5  & Scale separation
   & E & Global vacuum integrals
   & Bounce $\sim 10^{-32}$ of spacetime volume integrals
   & II \\
6  & Attractor-sensitivity dilemma
   & F & Bounce-set initial conditions
   & Attractors erase $\phi_i$; no-attractor needs $10^{-30}$ precision
   & II \\
7  & Parameter immunity
   & G & Cyclic vacuum selection
   & Action parameters $\neq$ boundary conditions;\; cyclic geometry incompatible
   & II \\
8  & Parity-even interaction
   & H & Tensor chirality
   & $(J^5)^2$ is parity-even;\; $\tilde{R}R = 0$ on FRW
   & III \\
9  & Phase-space conservation
   & J & Reversible state selection
   & Liouville theorem at $H=0$;\; no phase-space contraction
   & IV \\
10 & UV--IR specificity dilemma
   & L & Observable mechanism search
   & Generic mechanisms don't need the bounce
   & III \\
11 & Decoupling universality
   & L & PGT lower-scale bounce
   & Light torsion decouples: $g_\mathrm{eff} \sim m_T/M_\mathrm{Pl}^2$
   & I \\
12 & Vacuum amplification ceiling
   & M & Detectable GW from bounce
   & $\Omega_\mathrm{GW} \propto (H_b/M_\mathrm{Pl})^2$;\; gap $\geq 10^{17}$ to detectors
   & III \\
13 & Gravitational democracy
   & N & Relic/baryon production
   & Torsion $\sim 1/g_* \sim 1\%$ of Planck-era channels;\; $B$, $L$ conserved
   & V \\
14 & Bounce--vacuum decoupling
   & O & Irreversible vacuum selection
   & Trigger (bounce) $\neq$ outcome ($\rho_\mathrm{DE}$)
   & V \\
\hline\hline
\end{tabular}
\end{table*}


% ---------------------------------------------------------------------------
\subsection{Five Failure Modes}
\label{sec:failure_modes}

The fourteen barriers organize into five failure modes that together constitute
a comprehensive obstruction theorem for the minimal ECH framework.

\medskip
\noindent\textbf{Mode I: Geometric specificity lost on FRW
  (Barriers 1--4, 11).}
The ECH Lagrangian contains rich geometric structure---propagating torsion
modes, the Nieh--Yan topological density, non-metricity, disformal couplings.
On a Friedmann--Robertson--Walker background, however, isotropy and homogeneity
set $T_\mu = Q_\mu = 0$.  The only surviving geometric scalar is the Ricci
curvature $R(g)$, which is the same in general relativity and in all
metric-compatible theories.  Concretely:
\begin{itemize}
\item The PGT torsion mode's mass and coupling are both controlled by
  a single PGT parameter $t_3$, creating a lock: making the mass cosmologically
  light simultaneously decouples the mode from matter (Barrier~1).
\item The Nieh--Yan density $dT^a \wedge e_a$ is non-topological in MAG, but
  shift symmetry (which protects a small mass) and non-trivial geometric
  content are mutually exclusive for linear couplings (Barrier~2).
\item Any torsion- or non-metricity-descended scalar on FRW reduces to a
  standard scalar-tensor field $\phi$ with coupling $\xi\phi^2 R$---the same
  theory obtained from any non-minimal scalar, with no geometric fingerprint
  (Barrier~3).
\item Disformal couplings require two field gradients per vertex, but the
  connection coupling $\nabla = \partial + \Gamma$ provides only one,
  suppressing all distinctive effects to dimension-8 operators at
  $\mathcal{O}(M_\mathrm{Pl}^{-4})$ (Barrier~4).
\end{itemize}

\medskip
\noindent\textbf{Mode II: Scale separation (Barriers 5--7).}
The bounce occurs at $\rho_\mathrm{crit} \sim 0.27\,\rho_\mathrm{Pl}$, while
dark energy operates at $\rho_\mathrm{DE} \sim 10^{-122}\,\rho_\mathrm{Pl}$.
This $10^{122}$ hierarchy in energy density (equivalently $10^{61}$ in mass
scale) prevents the bounce from communicating with the DE sector:
\begin{itemize}
\item Global spacetime integrals are dominated by the expansion era; the bounce
  contributes a fraction $\sim 10^{-32}$ of the total spacetime volume
  (Barrier~5).
\item Quintessence initial conditions either flow to late-time attractors
  (erasing all bounce information) or require $10^{-30}$ precision in
  $\phi_i$, which the generic bounce cannot provide (Barrier~6).
\item Action parameters (coupling constants) cannot be constrained by dynamical
  boundary conditions, and cyclic cosmology with $\Lambda > 0$ and
  Planck-density bounce is geometrically incompatible with observed spatial
  flatness (Barrier~7).
\end{itemize}

\medskip
\noindent\textbf{Mode III: Generic or undetectable signatures
  (Barriers 8, 10, 12).}
Where the bounce does produce calculable perturbation spectra, the results are
either undetectable or indistinguishable from other bouncing cosmologies:
\begin{itemize}
\item The four-fermion interaction $(J^5)^2$ is parity-\emph{even}
  (pseudovector squared = scalar), and the Pontryagin density $\tilde{R}R$
  vanishes identically on FRW.  No parity-odd tensor signal exists
  (Barrier~8).
\item The tensor power spectrum $P_T \sim 2 \times 10^{-64}$ is calculable and
  scale-invariant ($n_T \approx 0$) but lies at least $10^4$ below the
  sensitivity of any planned detector.  The scalar transfer function
  $\mathcal{T}(k) = 1$ for $k \ll k_b$ by time-reversal symmetry, and
  features appear only at $k \sim k_b \sim \mathrm{GHz}$
  ($k_\mathrm{CMB}/k_b \sim 10^{-28}$).
\item Generic UV$\to$IR mechanisms (contraction spectra, inflation matching,
  phase transitions) operate identically in any hot cosmology and do not
  require or probe the bounce (Barrier~10).  The vacuum amplification ceiling
  $\Omega_\mathrm{GW} \propto (H_b/M_\mathrm{Pl})^2$ applies universally to
  all bounce scenarios (Barrier~12).
\end{itemize}

\medskip
\noindent\textbf{Mode IV: Hamiltonian conservation (Barrier 9).}
The bounce is a conservative scattering event at $H = 0$.  Liouville's theorem
guarantees that phase-space volume is preserved: the bounce \emph{rotates}
dark-sector states but cannot \emph{select} them.  Five candidate selection
mechanisms (pNGB misalignment, multi-vacuum, symmetry re-breaking, metastable
trapping, nonadiabatic excitation) all fail because the bounce provides no
dissipation.

\medskip
\noindent\textbf{Mode V: Gravitational democracy and vacuum decoupling
  (Barriers 13--14).}
At $T \sim M_\mathrm{Pl}$, the torsion interaction is one of approximately
$g_* \sim 100$ gravitational-strength channels, contributing $\sim 1\%$ of the
total scattering rate.  It cannot dominate relic production or baryogenesis,
and it conserves $B$ and $L$.  Meanwhile, irreversible processes (phase
transitions, tunneling) can be triggered by the bounce but the resulting vacuum
energy is set by hidden-sector parameters that the bounce does not determine.
The bounce controls \emph{whether} a transition occurs, not \emph{what} vacuum
energy results.


% ---------------------------------------------------------------------------
\subsection{Branch Status Summary}
\label{sec:branch_summary}

Table~\ref{tab:branches} summarizes the status of all fifteen branches.

\begin{table}[h]
\centering
\caption{Status of the fifteen branches investigated.  CLOSED: structural
  barrier identified.  GENERIC: result correct but shared by all bouncing
  cosmologies.  WEAK: no nontrivial constraint found.}
\label{tab:branches}
\begin{tabular}{cllc}
\hline\hline
Branch & Route tested & Barriers & Status \\
\hline
A & Propagating torsion (PGT) & 1 & Closed \\
B & Nieh--Yan pseudoscalar & 2 & Closed \\
C & Environmental mass & 3 & Closed \\
D & Disformal coupling & 4 & Closed \\
E & Global vacuum integrals & 5 & Closed \\
F & Bounce-set initial conditions & 6 & Closed \\
G & Cyclic vacuum selection & 7 & Closed \\
H & Tensor spectrum \& chirality & 8 & Closed \\
I & Horndeski DE stability & --- & Weak \\
J & State selection & 9 & Closed \\
K & Scalar perturbations & --- & Generic \\
L & UV--IR bridge (PGT) & 10, 11 & Closed \\
M & PGT bounce GW spectrum & 12 & Generic \\
N & Baryogenesis \& relics & 13 & Closed \\
O & Hidden-sector vacuum & 14 & Closed \\
\hline\hline
\end{tabular}
\end{table}


% ---------------------------------------------------------------------------
\subsection{Implications}
\label{sec:closure_implications}

The fourteen barriers organize into a coherent obstruction.  The bounce is
trapped between three regimes:

\begin{enumerate}
\item \emph{Too high-energy for late-time observables.}  The bounce energy
  density $\rho_\mathrm{crit} \sim 0.27\,\rho_\mathrm{Pl}$ is $10^{122}$
  times the dark energy scale.  All direct channels from bounce to DE are
  suppressed by this hierarchy (Modes I and II).

\item \emph{Too generic at its own scale.}  At $T \sim M_\mathrm{Pl}$,
  torsion is one of $\sim 100$ gravitational-strength interactions, and tensor
  and scalar spectra are shared by all symmetric radiation bounces (Modes III
  and V).

\item \emph{Too decoupled from late-time vacuum.}  Conservative dynamics
  (Liouville) preserve phase space, and irreversible processes decouple the
  trigger from the outcome (Mode IV and V).
\end{enumerate}

\noindent These results do not invalidate the ECH framework.  The bounce
remains a well-defined, singularity-free cosmology.  What they establish is
that the \emph{minimal} framework is \emph{observationally inert} in the
direct sector: no distinctive, testable consequence of the bounce dynamics
itself is accessible.  Dark energy requires separate input; we treat it as a
cosmological constant~$\Lambda$.

One \emph{indirect} handle survives.  The ECH parity-odd sector motivates a
Planck-scale axion-like particle whose photon coupling produces cosmic
birefringence.  This is not a bounce prediction---it is a prediction of the
effective ALP sector that the framework motivates.  We turn to this in the
next section.
